\subsection{Decreasing sequence of eigenvalues}

In this section, we assume that K = C.

We set \(\overline{\mathbf{N}} = \mathbf{N} \cup \{+\infty\} \subset \overline{\mathbf{R}}\) . In this section, we say that a vector space E has dimension \(+\infty \in \overline{\mathbf{N}}\) if E is not finite-dimensional. Let \(\mathrm{I_E} \subset \mathbf{N}\) be the set of dimensions of finite-dimensional subspaces F of E such that \(\mathrm{F} \neq \mathrm{E}\) . We have \(\mathrm{I_E} = \mathbf{N}\) if E is infinite-dimensional, and otherwise \(\mathrm{I_E} = \{0, \ldots, \dim(\mathrm{E}) - 1\}\) . We set \(\mathrm{I} = \mathrm{I_E}\) when no confusion can result.

Let \(u\) be a compact positive (in particular Hermitian) endomorphism of E. The sensitive spectrum of \(u\) is the set of values of a strictly decreasing sequence \((\nu_{k})_{0\leqslant k<\mathrm{Card}(\mathrm{Sp}_{\mathrm{s}}(u))}\) of positive real numbers (cf.

prop. 5 of III, p. 90). For every integer \(k\) such that \(0 \leqslant k < \text{Card}(\text{Sp}_s(u))\) , we denote by \(n_k\) the spectral multiplicity of \(\nu_k\), with \(n_k \geqslant 1\) . Let \(\text{M} \in \overline{\mathbf{N}}\) be the sum of the spectral multiplicities \(n_k\) ; it is the dimension of the image of \(u\) . We have \(\text{M} \leqslant \text{Card}(\text{I})\) .

For \(0 \leqslant n < M\) , one defines \(\lambda_n(u) = \nu_k\) , where \(k \geqslant 0\) is the unique integer such that

\[
n _ {0} + \dots + n _ {k - 1} \leqslant n <   n _ {0} + \dots + n _ {k}.
\]

One sets \(\lambda_{n}(u)=0\) if \(n\in I\) satisfies \(n\geqslant M\) . This case can occur only if I=N and if \(\mathrm{Sp}_{\mathrm{s}}(u)\) is finite (or, equivalently, if E is infinite-dimensional and u is of finite rank).

The sequence \((\lambda_{n}(u))_{n\in\mathrm{I}}\) is decreasing; for every \(\lambda\in\mathrm{Sp}_{\mathrm{s}}(u)\) , the number of integers n such that \(\lambda_{n}(u)=\lambda\) is equal to the spectral multiplicity of the eigenvalue \(\lambda\) of u.

One says, by abuse of language, that \((\lambda_{n}(u))_{n\in\mathrm{J}}\) is the decreasing sequence of eigenvalues of u repeated with their multiplicities.

PROPOSITION 3. --- Let u be a compact positive endomorphism of E. There exists an orthonormal family \((e_{n})_{n\in\mathrm{I}}\) in E such that, for every \(x\in E\) , one has

\[
u (x) = \sum_ {n \in \mathrm{I}} \lambda_ {n} (u) \langle e _ {n} | x \rangle e _ {n}, \quad \langle x | u (x) \rangle = \sum_ {n \in \mathrm{I}} \lambda_ {n} (u) | \langle e _ {n} | x \rangle | ^ {2}.
\]

Let \(\mathrm{B} = (f_{j})_{j \in \mathrm{J}}\) be an orthonormal basis of E in which u is diagonal (cor. 1 of IV, p. 150) and \((\lambda_{j})_{j \in \mathrm{J}}\) the family of eigenvalues of u in the basis B. Let J' be the set of \(j \in J\) such that \(\lambda_{j}\) belongs to the sensitive spectrum of u.

For each \(\lambda \in \mathrm{Sp}_{\mathrm{s}}(u)\) , there exists a bijection between the integers \(n\) such that \(0 \leqslant n < \mathrm{M}\) and \(\lambda_{n}(u) = \lambda\) and the \(j \in \mathrm{J}'\) such that \(\lambda_{j} = \lambda\) , since these two sets have the same cardinality, namely the spectral multiplicity of \(\lambda\) . A choice of such bijections for every \(\lambda\) defines a bijection \(\iota\) from the set of integers such that \(0 \leqslant n < \mathrm{M}\) to the set \(\mathrm{J}'\) . One defines the sequence \((e_n)_{0 \leqslant n < \mathrm{M}}\) in E by setting \(e_n = f_{\iota(n)}\) for \(0 \leqslant n < \mathrm{M}\) . It is an orthonormal family in E.

In the case where M \textless{} Card(I) = +∞, the space F generated by \(\{e_{0},\ldots,e_{M-1}\}\) is finite-dimensional and its orthogonal F° is infinite-dimensional; one chooses for \((e_{n})_{n\geqslant M}\) an orthonormal family in F°.

For every \(x \in \mathrm{E}\) , we have

\[
u (x) = \sum_ {j \in J ^ {\prime}} \lambda_ {j} \langle f _ {j} \mid x \rangle f _ {j} = \sum_ {0 \leqslant n <   M} \lambda_ {n} (u) \langle e _ {n} \mid x \rangle e _ {n}
\]

since \(u(f_{j}) = 0\) when \(j \in J - J'\) . If \(n \in I\) satisfies \(n \geqslant M\) , we have \(\lambda_{n}(u) = 0\) , and one obtains the first formula of the proposition. The second follows from it.

PROPOSITION 4. --- Let u be a compact positive endomorphism of E. Let \(f \in \mathcal{C}(\mathbf{R}_{+})\) be a continuous increasing function such that \(f(0) = 0\) .

The endomorphism \(f(u)\) is compact and positive and, for every \(n \in \mathrm{I}_{\mathrm{E}}\) , we have \(\lambda_n(f(u)) = f(\lambda_n(u))\) .

The endomorphism \(f(u)\) is compact and positive by prop. 6, b) of III, p. 91 and prop. 15, a) of I, p. 117. The spectrum of \(f(u)\) is the image under \(f\) of the spectrum of \(u\) (cor. 2 of I, p. 111). If \(\lambda \in \mathrm{Sp}_{\mathrm{s}}(f(u))\) , the spectral multiplicity of \(\lambda\) is the sum of the spectral multiplicities of the \(\mu \in \mathrm{Sp}(u)\) such that \(f(\mu) = \lambda\) (cor. 2 of III, p. 84). Since \(f\) is increasing, the sequence \((f(\lambda_n(u)))_{n \in \mathrm{I_E}}\) is decreasing. The assertion then follows from the definition of the sequence \((\lambda_n(f(u)))_{n \in \mathrm{I_E}}\) .

